\MakeAbstract{
    Unix V6++ 是同济大学操作系统课程使用的教学内核。该系统保留了 Unix V6 的基本结构，代码规模较小，便于展示进程管理、内存管理、文件系统和设备管理等核心机制。其原始实现面向 IA32 架构，并采用 C++ 编写；随着教学环境和系统软件技术的发展，该实现在线程安全边界、资源生命周期管理和体系结构适配方面逐渐暴露出维护成本。

    本课题以 Unix V6++ 的内存管理和文件系统为主要对象，分析原系统中页分配、交换管理、inode 缓存、打开文件表和缓冲区缓存等机制的实现方式，并在保持教学结构和 Unix V6 风格设计的基础上使用 Rust 进行重构。重构内容包括类型化地址抽象、Buddy 与 Slab 分层内存分配、引用封装、RAII 风格资源回收、生命周期约束和显式错误传播。课题同时完成了面向 RISC-V 64 位平台的底层适配，涉及 Sv39 分页、MMIO 串口、VirtIO 块设备和 QEMU 运行环境。

    测试结果表明，Rust + IA32 和 Rust + RISC-V 两个版本均能够完成文件系统装载、shell 运行、文件读写和用户程序执行等基本任务。系统调用、内存分配和 I/O 吞吐量测试显示，重构后的文件系统和内存管理路径具备可接受的运行开销。相较于原始 C++ 实现，Rust 版本在对象生命周期、错误返回和模块边界方面表达更明确，为 Unix V6++ 的后续教学使用和体系结构迁移提供了基础。
}{Unix V6++，Rust，RISC-V，内存管理，文件系统，操作系统教学}

\MakeAbstractEng{
    Unix V6++ is a teaching operating system kernel used in the operating system course at Tongji University. It preserves the basic structure of Unix V6 and keeps the code base small enough for teaching process management, memory management, file systems, and device management. The original implementation targets the IA32 architecture and is written in C++. As the teaching environment and system software practice evolve, this implementation shows increasing maintenance costs in resource lifetime management, safety boundaries, and architecture adaptation.

    This thesis focuses on the memory management and file system of Unix V6++. It analyzes page allocation, swap management, inode caching, open file tables, and buffer cache mechanisms in the original system, and then refactors the key modules in Rust while preserving the teaching-oriented structure and Unix V6 style design. The refactoring introduces typed address abstractions, a Buddy and Slab based layered allocator, reference wrappers, RAII-style resource reclamation, lifetime constraints, and explicit error propagation. The thesis also adapts the system to the RISC-V 64-bit platform, including Sv39 paging, MMIO UART, VirtIO block devices, and the QEMU execution environment.

    The test results show that both Rust + IA32 and Rust + RISC-V versions can mount the file system, run the shell, perform file operations, and execute user programs. System call, memory allocation, and I/O throughput tests indicate that the refactored memory management and file system paths have acceptable runtime overhead. Compared with the original C++ implementation, the Rust version expresses object lifetimes, error returns, and module boundaries more explicitly, providing a basis for further teaching use and architecture migration of Unix V6++.
}{Unix V6++, Rust, RISC-V, memory management, file system, operating system education}
